\begin{thebibliography}{00}
	
	\makeatletter
	\renewcommand\@biblabel[1]{#1}
	\makeatother	
	\bibitem[]{} \textbf{{\large Нормативные акты:}}
	\makeatletter
	\renewcommand\@biblabel[1]{#1.}
	\makeatother
	
	\bibitem{book1} Об искусственном интеллекте [Электронный ресурс] : [закон Республики Казахстан от 17 ноября 2025 года № 230-VIII : по состоянию на 6 марта 2026 г.]. – Доступ из информационной системы «Параграф».
	
	
%	\bibitem{book2} Python для сложных задач: наука о данных и машинное обучение / Плас Дж. Вандер : пер. с англ. -- СПб. :  ООО «Питер», 2018. С. 576. URL: \url{https://iro23.ru/sites/default/files/plas_dzh._vander_-_python_dlya_slozhnyh_zadach_nauka_o_dannyh_i_mashinnoe_obuchenie_bestsellery_oreilly_-_2018.pdf} (дата обращения: 28.01.2025).
%	\bibitem{book3} Python и машинное обучение: машинное и глубокое обучение с использованием Python, scikit-learn и TensorFlow 2 / С. Рашка, В. Мирджалили : под редакцией Ю. Н. Артеменко. 3-е изд., пер. с англ. -- СПб. : ООО «Диалектика», 2020. С. 848. URL: \url{https://coollib.net/b/674877-sebastyan-rashka-python-i-mashinnoe-obuchenie} (дата обращения: 28.01.2025).
	
	\makeatletter
	\renewcommand\@biblabel[1]{#1}
	\makeatother	
	\bibitem[]{} \textbf{{\large Статьи:}}
	\makeatletter
	\renewcommand\@biblabel[1]{#1.}
	\makeatother	
	
	\bibitem{article1} Бобкова, Е.В. Применение искусственного интеллекта в банковской сфере: технологии и результаты // Выпускная квалификационная работа Санкт-Петербургский государственный университет. -- 2025. -- С. 1 -- 72.
	
	\bibitem{article2} Гозгешев, Э.А. Искусственный интеллект в банковском секторе: возможности и риски // Вестник евразийской науки. – 2024. – № 16, том 6. – С. 1 – 9.
	
	\bibitem{article3} Городецкая, О.Ю. Ключевые тренды применения искусственного интеллекта в банковской сфере / О.Ю. Городецкая, Я.Л. Гобарева // Финансовые рынки и банки. -- 2022. -- № 12. -- C. 34 -- 42.
	
	\bibitem{article4} Зиятбекова, Г.З. Исследование систем кредитного скоринга с использованием алгоритмов машинного обучения / Г.З. Зиятбекова, Д.К. Даркенбаев, Д.К. Алтыбай, А. Даркенбаев, Н.О. Мекебаев // КазУТБ. -- 2024. -- № 3, том 24. -- С. 1 -- 11.
	
	\bibitem{article5} Касенова, А.А. Определение кредитного скоринга заемщиков банка на основе технологии машинного обучения // Курсовая работа. -- 2025 -- С. 1 -- 27.
	
	\bibitem{articleLemeshko1} Лемешко, Б.Ю. О распределениях статистик и мощности критериев однородности двух выборок / Б.Ю. Лемешко, С.Б. Лемешко // Материалы Российской НТК “Информатика и проблемы коммуникаций”. -- 2005. -- Т. 1. -- С. 109--112.
	
	\bibitem{articleLemeshko2} Лемешко, Б.Ю. Исследование распределения оценок коэффициента корреляции в зависимости от истинного значения корреляции / А.В. Танасейчук, Б.Ю. Лемешко // Материалы VIII международной конференции “Актуальные проблемы электронного приборострое­ния”. -- 2006. -- Т. 6. -- С. 91--94.
	
	\bibitem{article6} Михайлов, Д.И. Применение методов искусственного интеллекта в системе риск-менеджмента банка // Научные записки молодых исследователей. -- 2025. -- № 13, том 3. -- С. 55 -- 61.
	
	\bibitem{article7} Смирнов, С.В. Методы машинного обучения в макроэкономическом прогнозировании: предварительные итоги // Вопросы экономики. -- 2025. -- № 10. -- С. 131 -- 154.
	
	\bibitem{articleSmirnov} Смирнов, Н.В. Оценка расхождения между эмпирическими кривыми распределения в двух независимых выборках // Бюллетень МГУ, серия А. -- 1939. -- №2, том 2. -- С. 3--14.
	
	\bibitem{article8} Addy, W.A. AI in credit scoring: A comprehensive review of models and predictive analytics / W.A.Addy, A.O. Ajayi-Nifise, B.G. Bello, S.T. Tula, O. Odeyemi, T. Falaiye // Global Journal of Engineering and Technology Advances. -- 2024. -- № 18, том 2. -- С. 118 -- 129.
	
	
	\bibitem{article9} Bary, H. A systematic literature review of models and predictive analytics / H. Bary, S. Juthi, A.M. Mistry, Kamrujjaman // Journal of machine learning, data engineering and data science. -- 2024. -- № 1, том 1. -- С. 1 -- 18. 
	
	\bibitem{Fisher9} Fisher, R.A. Frequency Distribution of the Values of the Correlation Coefficient in Samples from an Indefinitely Large Population // Biometrika. -- 1915. -- №4, том 10. -- С. 507--521.
	
	\bibitem{article10} Kothandapani, H.P. Leveraging AI for credit scoring and financial inclusion in emerging markets // World Journal of Advanced Research and Reviews. -- 2022. -- № 15, том 3. -- С. 526 -- 539.
	
	\bibitem{article11} Mehndiratta, N. The use of Artificial Intelligence in the Banking Industry / N. Mehndiratta, G. Arora, R. Bathla // International Conference on Recent Advances in Electrical, Electronics \& Digital Healthcare Technologies (REEDCON). -- 2023. -- C. 588 -- 591.
	
	\bibitem{article12} Nuka, T.F. AI and Machine learning as tools for financial inclusion: challenges and opportunities in credit scoring // International Journal of Science and Research Archive. -- 2026. -- № 13, том 2. -- С. 1052 -- 1067.
	
	\bibitem{article13} Omotosho, M.A. Exploring the role of AI-driven credit scoring systems on financial inclusion in emerging economies // American Journal of Humanities and Social Sciences Research. -- 2025. -- № 9, том 2. -- С. 1 -- 10.
	
	\bibitem{article14} Vakrani, D.S. Evaluating AI-driven credit scoring models versus traditional statistical techniques / D.S. Vakrani, P.S. Padhye, S.K. Gupta, J.K. Roy, Nerlekar V.S., Parashar N. // Discover Artificial Intelligence. --2026. -- № 6. -- С. 1 -- 14.
	
	
	\bibitem{article15} Wang, H. Application of Decision Tree Model in Personal Credit Scoring and Its Fairness Optimization // Advances in Economics, Management and Political Sciences. -- 2025. -- С. 109--118.
	
	
	\makeatletter
	\renewcommand\@biblabel[1]{#1}
	\makeatother
	\bibitem[]{} \textbf{{\large Интернет-ресурсы:}}
	\makeatletter
	\renewcommand\@biblabel[1]{#1.}
	\makeatother
	
	
	\bibitem{kursiv1}Забытые долги казахстанцев перед банками выросли в полтора раза за год [Электронный ресурс] / Общественно-политическое издание kursiv.kz -- URL: \url{https://kz.kursiv.media/2026-03-02/fvfv-zabytyye-dolgi-kazakhstantsev-pered-bankami-vyrosli-v-poltora-raza-za-god/}. -- Загл. с экрана.
	
	\bibitem{nbrb} Искусственный интеллект и его применение в банковском секторе [Электронный ресурс] / Национальный банк Республики Беларусь. -- URL: \url{https://www.nbrb.by/bv/pdf/articles/12217.pdf}. -- Загл. с экрана.
	
	\bibitem{nbrk} Искусственный интеллект на финансовых рынках Центральной Азии. Текущее состояние и перспективы [Электронный ресурс] / Национальный банк Республики Казахстан. -- URL: \url{https://nationalbank.kz/ru/news/analitika-metodologiya/rubrics/2211}. -- Загл. с экрана.
	
	\bibitem{kursiv2} Казахстанцев ждет долговая усталость в 2026 году: экономист озвучил опасения [Электронный ресурс] / Общественно-политическое издание kursiv.kz -- URL: \url{https://kz.kursiv.media/2026-01-27/svan-kazahstancy-prodolzhayut-popolnyat-depozity-i-brat-novye-kredity-nesmotrya-na-stavku/}. -- Загл. с экрана.
	
	\bibitem{cninform} Казахстан занимает 76-е место в списке стран с наибольшей готовностью к развитию искусственного интеллекта [Электронный ресурс] / Казахстанское международное информационное агентство. -- URL: \url{https://cn.inform.kz/news/zairengongnengfazhongdi76-fed5aa/}. -- Загл. с экрана.
	
	\bibitem{invento} Как использовать машинное обучение для принятия решений о кредитовании?  [Электронный ресурс] / Аутсорсинговая компания Invento Labs. -- URL: \url{https://invento-labs.by/o-kompanii/blog/2024/jan/kak-ispolzovat-mashinnoe-obuchenie-dlya-prinyatiya-reshenij-o-kreditovanii.html}. -- Загл. с экрана.
	
	\bibitem{yandexPCA} Как работает метод главных компонент (PCA) -- на простом примере  [Электронный ресурс] / Яндекс Практикум. -- URL: \url{https://practicum.yandex.ru/blog/metod-glavnyh-komponent/}. -- Загл. с экрана.
	
	\bibitem{zakon.kz} Ловушка в 20\%: почему казахстанцы продолжают кредитоваться при падающих доходах [Электронный ресурс] / Информационный портал РК. -- URL: \url{https://www.zakon.kz/finansy/6505800-lovushka-v-20-pochemu-kazakhstantsy-prodolzhayut-kreditovatsya-pri-padayushchikh-dokhodakh.html}. -- Загл. с экрана.
	
	\bibitem{Tengrinews} Можно ли «отбелить» кредитную историю и есть ли черный список должников в Казахстане [Электронный ресурс] / Информационный интернет-портал «Tengrinews». -- URL: \url{https://tengrinews.kz/kazakhstan_news/li-otbelit-kreditnuyu-istoriyu-est-chernyiy-spisok-doljnikov-560869/}. -- Загл. с экрана.
	
	\bibitem{ZestAI} Платформа оценки кредитоспособности на основе искусственного интеллекта «Zest AI» меняет будущее кредитования [Электронный ресурс] / Японская платформа облачного кредитования «CrowdLoan». -- URL: \url{https://corp.crowdloan.jp/axlab/unf-589-waid/}. -- Загл. с экрана.
	
	\bibitem{Emagia} Полное руководство по автоматизированному кредитному скорингу: глубокое погружение в ИИ, машинное обучения и будущее финансовых решений [Электронный ресурс] / Компания в области интеллектуального управления дебиторской задолженностью на основе искусственного интеллекта «Emagia». -- URL: \url{https://www.emagia.com/ru/resources/glossary/automated-credit-scoring/}. -- Загл. с экрана.
	
	\bibitem{24kz} Свыше 9 млн казахстанцев имеют кредиты на сумму больше 24 трлн тенге [Электронный ресурс] / Информационный портал 24.kz. -- URL: \url{https://24.kz/ru/news/social/729252-svyshe-9-mln-kazakhstantsev-imeyut-kredity-na-summu-bolshe-kzt24-trln}. -- Загл. с экрана.
	
	\bibitem{AlfaBank} Соревнование на данных кредитный историй [Электронный ресурс] / Open Data Science. -- URL: \url{https://ods.ai/competitions/dl-fintech-bki}. -- Загл. с экрана.
	
	\bibitem{Oxford} Government AI Readiness Index 2025 [Электронный ресурс] / Британская аналитическая и консалтинговая компания «Oxford Insights». -- URL: \url{https://oxfordinsights.com/wp-content/uploads/2026/01/2025-Government-AI-Readiness-Index-Report_01_26.pdf}. -- Загл. с экрана.
	
	\bibitem{GridSearchCV} GridSearchCV [Электронный ресурс] / Python-библиотека для машинного обучения. -- URL: \url{https://scikit-learn.org/stable/modules/generated/sklearn.model_selection.GridSearchCV.html}. --Загл. с экрана.
	
	\bibitem{PCA} PCA [Электронный ресурс] / Python-библиотека для машинного обучения. -- URL: \url{https://scikit-learn.org/stable/modules/generated/sklearn.decomposition.PCA.html}. -- Загл. с экрана.

	\bibitem{RandomizedSearchCV} RandomizedSearchCV [Электронный ресурс] / Python-библиотека для машинного обучения. -- URL: \url{https://scikit-learn.org/stable/modules/generated/sklearn.model_selection.RandomizedSearchCV.html}. -- Загл. с экрана.
	
\end{thebibliography}